\documentclass[11pt,a4paper]{article}

\usepackage[margin=1in]{geometry}
\usepackage[T1]{fontenc}
\usepackage{lmodern}
\usepackage{microtype}
\usepackage{enumitem}
\usepackage{xcolor}
\usepackage{hyperref}
\usepackage{listings}

\hypersetup{
    colorlinks=true,
    linkcolor=black,
    urlcolor=blue
}

\lstdefinestyle{commandstyle}{
    basicstyle=\ttfamily\small,
    backgroundcolor=\color{gray!8},
    frame=single,
    breaklines=true,
    columns=fullflexible
}

\setlist[itemize]{leftmargin=*}
\setlist[enumerate]{leftmargin=*}
\emergencystretch=2em

\title{FlexQL Design Document}
\author{}
\date{}

\begin{document}

\maketitle

\section{Repository}

\url{https://github.com/Somak-2001/FlexQL.git}

\section{Overview}

FlexQL is a lightweight SQL-like client/server database driver implemented fully in C++. The system is split into:

\begin{itemize}
    \item A TCP database server that parses and executes SQL-like commands
    \item A C-style client API layer that communicates with the server
    \item An interactive REPL built on top of the client API
\end{itemize}

The implementation supports \texttt{CREATE TABLE}, \texttt{INSERT}, \texttt{SELECT}, single-condition \texttt{WHERE}, and \texttt{INNER JOIN}.

\section{Storage Design}

\subsection{Row-Major Storage}

The database uses a row-major in-memory layout:

\begin{itemize}
    \item Each table stores a vector of rows
    \item Each row stores a vector of string values aligned with the schema
    \item This makes inserts simple and keeps row reconstruction cheap during \texttt{SELECT *}
\end{itemize}

\subsection{Durable Storage}

FlexQL persists data in a small WAL-plus-snapshot layout under \texttt{flexql\_data/}.

\begin{itemize}
    \item \texttt{catalog.txt} stores table schemas
    \item \texttt{tables/<TABLE>.rows} stores checkpointed live rows
    \item \texttt{wal.log} stores mutating SQL statements that have been synced but not yet checkpointed
\end{itemize}

The write path is:

\begin{enumerate}
    \item Parse and validate the statement
    \item Append the normalized mutating SQL statement to the WAL and \texttt{fsync} it
    \item Apply the mutation in memory
    \item Mark the affected table or catalog as dirty
    \item Rewrite dirty catalog/table snapshots atomically at checkpoint time
    \item Clear the WAL only after checkpoint completion
\end{enumerate}

On restart, the engine loads the catalog and table snapshots first, then replays any remaining WAL entries. This means a crash after the WAL append but before checkpointing is recoverable: the snapshot may be older, but the synced WAL entry is replayed to reconstruct the committed state.

Disk is the durable source of truth. RAM is used only as the active working set: table rows, primary-key hash indexes, and cached \texttt{SELECT} results are loaded and maintained in memory so query execution avoids disk reads on the hot path. This improves assignment-scale latency and keeps the implementation simple, but it means very large tables require enough RAM for their active rows and indexes. The snapshots and WAL make the in-memory state reconstructible after restart, while the LRU query cache remains an optimization only and can be discarded at any time.

The main persistence trade-off is checkpoint frequency. Checkpointing after every row would minimize WAL replay work but would add heavy disk I/O. FlexQL instead syncs each mutating statement to the WAL and checkpoints dirty snapshots in batches, which reduces write amplification while preserving recoverability.

\subsection{Schema Representation}

Each table keeps:

\begin{itemize}
    \item A list of columns with name, type, and primary-key metadata
    \item A hash map from column name to column index
    \item A row container
\end{itemize}

Supported types:

\begin{itemize}
    \item \texttt{DECIMAL}
    \item \texttt{VARCHAR}
\end{itemize}

Type checks are performed during insertion. Expiration timestamps are stored separately as absolute epoch seconds so TTL handling remains independent of the table's supported column types.

\section{Indexing Method}

Primary indexing is implemented as a hash map:

\begin{itemize}
    \item Key: primary-key value
    \item Value: row offset inside the table storage
\end{itemize}

This design makes exact-match lookups on the indexed column efficient and keeps insert overhead low.

Current assumption:

\begin{itemize}
    \item At most one primary-key column per table
\end{itemize}

\section{Query Execution}

\subsection{CREATE TABLE}

\begin{itemize}
    \item Parses schema metadata
    \item Verifies duplicate column names are not present
    \item Registers the table in the in-memory catalog
\end{itemize}

\subsection{INSERT}

\begin{itemize}
    \item Validates every row in the statement against schema width and type rules
    \item Enforces unique primary-key values against both existing rows and rows inside the same batch
    \item Supports multi-row \texttt{INSERT INTO ... VALUES (...), (...), ...}
    \item Stores an optional absolute expiration timestamp for each inserted row
\end{itemize}

\subsection{SELECT}

\begin{itemize}
    \item Supports \texttt{SELECT *} and projected columns
    \item Supports a single \texttt{WHERE} predicate
    \item Supports \texttt{INNER JOIN ... ON ...}
\end{itemize}

\subsection{WHERE Restriction}

\begin{itemize}
    \item Only one condition is supported
    \item No \texttt{AND} or \texttt{OR}
\end{itemize}

\section{Expiration Handling}

Each row stores an optional expiration timestamp in epoch seconds.

Supported insertion forms:

\begin{itemize}
    \item \texttt{TTL <seconds>}
    \item \texttt{EXPIRES AT 'YYYY-MM-DD HH:MM:SS'}
\end{itemize}

Expired rows are removed lazily during future access to the table. This avoids doing background cleanup work while still preserving correctness.

\section{Caching Strategy}

FlexQL uses an LRU result cache for \texttt{SELECT} queries.

\begin{itemize}
    \item Cache key: normalized SQL text
    \item Cache value: complete query result including columns and rows
    \item Capacity: 128 query results
\end{itemize}

The cache is invalidated whenever:

\begin{itemize}
    \item A table is created
    \item A row is inserted
    \item Expired rows are purged
\end{itemize}

This keeps cached results correct after writes.

\section{Multithreading Design}

The server follows a thread-per-client model:

\begin{itemize}
    \item The listening thread accepts incoming TCP connections
    \item Each client connection is handled by a worker thread
    \item Shared database state is protected with \texttt{std::shared\_mutex}
    \item The query cache has its own mutex
\end{itemize}

This allows multiple clients to issue queries concurrently while preventing race conditions on shared structures.

\section{Networking Design}

Client and server communicate through a simple framed protocol:

\begin{itemize}
    \item Every message is prefixed with a 32-bit payload length
    \item The client sends one SQL statement per request
    \item The server replies with result headers, row frames, and a terminal \texttt{END} frame
\end{itemize}

This avoids ambiguity from newline-based parsing and makes binary-safe message transfer straightforward.

\section{Performance Notes}

The implementation is designed to scale reasonably for assignment workloads through:

\begin{itemize}
    \item Row-major storage for cheap appends
    \item Hash-based primary indexing
    \item LRU caching for repeated reads
    \item Batch insert parsing to reduce client/server round trips
    \item WAL-first persistence with batched checkpoints to reduce disk write amplification
    \item Lazy expiration cleanup
\end{itemize}

The included benchmark utility can be used to measure:

\begin{itemize}
    \item Insert throughput
    \item Point-select latency
\end{itemize}

Local performance results measured on April 7, 2026:

\begin{itemize}
    \item TCP unit test: \texttt{./benchmark\_bin --unit-test} passed 21/21 against \texttt{./server} on \texttt{127.0.0.1:9000}
    \item Benchmark wrapper, 1M rows: \texttt{./benchmark 1000000} inserted 1,000,000 rows in 0.083905 sec, or 11,918,240 rows/sec; budget check passed against the 2.0 sec target
    \item Benchmark wrapper, 10M rows: \texttt{./benchmark 10000000} inserted 10,000,000 rows in 0.836931 sec, or 11,948,416 rows/sec; budget check passed against the 20.0 sec target
    \item Read-heavy check outside the insert benchmark: 100 repeated \path|SELECT * FROM CACHE_BENCH_20260407 WHERE ID = 42;| requests over the TCP REPL completed in 8.15 sec end-to-end, about 81.5 ms/request including client, protocol, and table-printing overhead
    \item Concurrent-client check outside the insert benchmark: 4 REPL clients issuing 25 cached point selects each completed 100 total selects in 1.98 sec, with no client error output
\end{itemize}

\section{Compilation Instructions}

\begin{lstlisting}[style=commandstyle,language=bash]
make
\end{lstlisting}

\section{Execution Instructions}

Benchmark workflow, Terminal 1:

\begin{lstlisting}[style=commandstyle,language=bash]
sh compile.sh
./server
\end{lstlisting}

Benchmark workflow, Terminal 2:

\begin{lstlisting}[style=commandstyle,language=bash]
./benchmark --unit-test
./benchmark 1000
\end{lstlisting}

Manual REPL workflow, Terminal 1:

\begin{lstlisting}[style=commandstyle,language=bash]
make flexql_repl
./server
\end{lstlisting}

Manual REPL workflow, Terminal 2:

\begin{lstlisting}[style=commandstyle,language=bash]
./flexql_repl 127.0.0.1 9000
\end{lstlisting}

Exit the REPL with:

\begin{lstlisting}[style=commandstyle]
.exit
\end{lstlisting}

\section{Verification}

Local verification completed for:

\begin{itemize}
    \item Build success with \texttt{make}
    \item SQL parser and engine smoke test with batch \texttt{INSERT}, \texttt{SELECT}, \texttt{WHERE}, \texttt{TTL}, and \texttt{INNER JOIN}
    \item Restart persistence test using a fresh engine instance after on-disk checkpointing
    \item TCP client/server round-trip validation with \texttt{./benchmark\_bin --unit-test}: 21/21 tests passed over \texttt{127.0.0.1:9000}
\end{itemize}

\end{document}
